\clearpage
\thispagestyle{empty}

\begingroup
\setlength{\parindent}{0pt}
\sffamily

\begin{tikzpicture}[remember picture,overlay]

% Dunkelblauer Kopfbereich
\fill[BookBlueDark]
  (current page.north west)
  rectangle
  ([yshift=-108mm]current page.north east);

% Blauer Trennstreifen
\fill[BookBlue]
  ([yshift=-108mm]current page.north west)
  rectangle
  ([yshift=-112mm]current page.north east);

% Überschrift
\node[
  anchor=north west,
  text=white,
  text width=150mm,
  align=left,
  inner sep=0pt
] at ([xshift=30mm,yshift=-34mm]current page.north west) {%
  {\fontsize{28}{34}\selectfont\bfseries
   About this book\par}

  \vspace{11mm}

  {\fontsize{13.5}{19}\selectfont
   A practical introduction for everyone interested in the
   simulation of high-voltage electrical power networks.\par}
};

% Haupttext
\node[
  anchor=north west,
  text=BookGray,
  text width=150mm,
  align=left,
  inner sep=0pt
] at ([xshift=30mm,yshift=-129mm]current page.north west) {%
  \begin{minipage}{150mm}
  \fontsize{9.5}{13.2}\selectfont

  {\color{BookBlueDark}
   \fontsize{10.4}{14.4}\selectfont
   Over the past 25 years, I have developed two educational
   toolboxes for LTspice. The first supports nonlinear load-flow
   calculations; the second is designed for the transient
   simulation of electrical power networks.\par}

  \vspace{2.5mm}

  The toolboxes contain the essential components required to build
  realistic power-system models:

  \begin{itemize}
    \setlength{\itemsep}{0pt}
    \setlength{\parsep}{0pt}
    \setlength{\topsep}{1mm}
    \setlength{\partopsep}{0pt}

    \item buses and sources;
    \item synchronous machines and transformers;
    \item overhead lines and power cables;
    \item loads and shunt reactors; and
    \item surge-protection devices, Petersen coils, and many
          other elements.
  \end{itemize}

  Some elements provide parameters that can be changed directly by
  the user. Others are supplied with carefully selected presets,
  which can also be overwridden whenever a study requires different
  values.

  \vspace{1.5mm}

  The transformer library ranges from simple models to sophisticated
  representations with a nonlinear iron core, hysteresis, and
  remanent flux. Eight vector groups and voltage levels from
  \SI{400}{\kilo\volt} down to \SI{400}{\volt} are included.
  Additional vector groups can be created by editing the
  macro-models, because all \texttt{.lib} files are freely accessible
  and editable.

  \vspace{1.5mm}

  Hundreds of examples for both toolboxes are available on my
  website, ranging from simple introductory circuits to complex
  networks. Many schematic files contain their own documentation,
  and the first nine examples are explained in detail in this book.

  \vspace{1.5mm}

  These toolboxes are intended especially for students of electrical
  engineering, but they are equally useful to teachers, engineers,
  and anyone who wants to explore power-system behaviour through
  simulation.

  \vspace{2.5mm}

  {\color{BookBlueDark}
   \fontsize{10.4}{14.4}\selectfont
   I hope you enjoy experimenting with these simulations as much as
   I do. Your feedback is always welcome.\par}
  \end{minipage}
};

% Trennlinie im Fußbereich
\draw[LightBlue,line width=0.8pt]
  ([xshift=30mm,yshift=34mm]current page.south west) --
  ([xshift=-30mm,yshift=34mm]current page.south east);

% Autor
\node[
  anchor=south west,
  text=BookBlueDark,
  inner sep=0pt,
  font=\fontsize{10.5}{15}\selectfont
] at ([xshift=30mm,yshift=23mm]current page.south west)
  {Friedrich Schmidt};

% Website
\node[
  anchor=south east,
  text=BookBlue,
  inner sep=0pt,
  font=\fontsize{10.5}{15}\selectfont\bfseries
] at ([xshift=-30mm,yshift=23mm]current page.south east)
  {www.friedrichschmidt.at};

\end{tikzpicture}

\null
\clearpage
\endgroup